\section{Probability as area under a density}

A density becomes useful only when it is integrated. The integral collects the
small contributions of density over an interval. Geometrically, this integral is
an area under the density curve.

If \(X\) has density \(f_X\), then
\[
P(a\leq X\leq b)=\int_a^b f_X(x)\,dx.
\]
This formula is the continuous counterpart of adding probability masses in the
discrete case.

\subsection*{Sums and integrals}

For a discrete random variable, probabilities are obtained by summing masses:
\[
P(X\in A)=\sum_{x\in A}p_X(x).
\]
For a continuous random variable with density, probabilities are obtained by
integrating density:
\[
P(X\in A)=\int_A f_X(x)\,dx.
\]
The two formulas have the same purpose: both accumulate probability over a set of
values. The difference is the way probability is distributed.

\begin{center}
\begin{tabular}{lll}
\toprule
Case & Local description & Accumulation rule \\
\midrule
Discrete & mass \(p_X(x)\) at points & sum masses \\
Continuous & density \(f_X(x)\) over intervals & integrate density \\
\bottomrule
\end{tabular}
\end{center}

\subsection*{Small intervals and local approximation}

Suppose \(f_X\) is continuous near a point \(x\). For a small positive number
\(h\), the probability of the interval \([x,x+h]\) is
\[
P(x\leq X\leq x+h)=\int_x^{x+h}f_X(t)\,dt.
\]
If the interval is very short, the density does not change much across it, so
\[
\int_x^{x+h}f_X(t)\,dt\approx f_X(x)h.
\]
This approximation explains why density can be read as local concentration of
probability. A larger density means that probability accumulates faster per unit
length near that point.

The approximation must not be confused with an equality in general. The exact
probability is the integral. The product \(f_X(x)h\) is a useful local
approximation when \(h\) is small and the density varies slowly over the
interval.

\subsection*{Example: constant density}

Let \(X\) be uniformly distributed on \([0,10]\). Then
\[
f_X(x)=
\begin{cases}
\frac1{10}, & 0\leq x\leq 10,\\
0, & \text{otherwise}.
\end{cases}
\]
For the interval \([2,7]\),
\[
P(2\leq X\leq 7)=\int_2^7 \frac1{10}\,dx=\frac{5}{10}=0.5.
\]
The probability is the area of a rectangle whose width is \(5\) and whose height
is \(1/10\).

For the shorter interval \([2,3]\),
\[
P(2\leq X\leq 3)=\int_2^3 \frac1{10}\,dx=0.1.
\]
In a constant-density model, probability is proportional to interval length.

\subsection*{Example: increasing density}

Let
\[
f_X(x)=
\begin{cases}
2x, & 0\leq x\leq 1,\\
0, & \text{otherwise}.
\end{cases}
\]
The probability of the first half of the interval is
\[
P(0\leq X\leq 0.5)=\int_0^{0.5}2x\,dx=0.25.
\]
The probability of the second half is
\[
P(0.5\leq X\leq 1)=\int_{0.5}^{1}2x\,dx=0.75.
\]
The two intervals have the same length, but the second interval receives three
times as much probability. The reason is not that it contains more points; both
intervals contain infinitely many points. The reason is that the density is
larger on the second interval.

\subsection*{Events formed from intervals}

If an event is a union of disjoint intervals, its probability is the sum of the
corresponding areas. For example, if \(X\) is uniformly distributed on \([0,1]\),
then
\[
P(X\in[0.1,0.2]\cup[0.7,0.9])
=
P(0.1\leq X\leq 0.2)+P(0.7\leq X\leq 0.9).
\]
The intervals are disjoint, so
\[
P(X\in[0.1,0.2]\cup[0.7,0.9])=0.1+0.2=0.3.
\]
This is the same additivity principle used in earlier chapters. What changes is
how the probability of each interval is computed.

\subsection*{Why single points have probability zero}

For a continuous random variable with density \(f_X\), the probability of a
single point \(a\) is
\[
P(X=a)=P(a\leq X\leq a)=\int_a^a f_X(x)\,dx=0.
\]
The interval has zero width, so the area is zero.

This argument is precise for random variables described by densities. It is not
a statement about every possible random variable. Some random variables may have
both continuous parts and point masses. In this chapter, unless stated
otherwise, continuous random variables are treated through densities.

\begin{remarkbox}
For a density model, changing or removing finitely many individual points does
not change the probability of an interval. Area is not affected by individual
points.
\end{remarkbox}

\subsection*{Checking probabilities by geometry}

A useful way to check continuous probability calculations is to compare the area
with the total area. Since the total area under a density is \(1\), an interval
cannot have probability larger than \(1\). If a computed interval probability is
negative or larger than \(1\), something is wrong.

Also, if one interval lies inside another, its probability cannot be larger:
\[
[a,b]\subseteq[c,d]
\quad\Longrightarrow\quad
P(a\leq X\leq b)\leq P(c\leq X\leq d).
\]
This monotonicity comes from the non-negativity of the density.

\begin{summarybox}
A density describes how probability is spread. An integral accumulates that
spread over an interval. The probability is the area, not the height.
\end{summarybox}
